\section{Limpieza y preprocesamiento}
\label{sec:preprocesamiento}

El preprocesamiento tuvo como propósito construir una base analítica
consistente sin alterar innecesariamente la información original. El
procedimiento se diseñó para ser reproducible, de modo que el notebook
pueda ejecutarse desde el inicio tanto si el dataset ya se encuentra en
el repositorio como si debe descargarse desde la fuente oficial.

A diferencia de otros problemas aplicados, el conjunto de datos no
requirió imputación, eliminación de registros ni corrección manual de
categorías. El trabajo se concentró en verificar la integridad de la
fuente, documentar la codificación de la respuesta, conservar la
naturaleza estadística de los predictores y construir las variables
derivadas necesarias para el análisis.

\subsection{Configuración reproducible}

Al inicio del notebook se fija una semilla global:

\[
\texttt{SEED}=2026.
\]

La semilla se utiliza en los procedimientos que dependen de generación
seudoaleatoria, entre ellos el remuestreo bootstrap, la separación de
datos, la validación cruzada y el sampler Metropolis--Hastings. Su uso
permite reproducir las principales estimaciones obtenidas durante el
análisis.

El entorno utilizado corresponde a Python 3.12.10. Las versiones
registradas durante la ejecución se presentan en la
Tabla~\ref{tab:entorno-computacional}.

\begin{table}[H]
    \centering
    \caption{Entorno computacional registrado por el notebook.}
    \label{tab:entorno-computacional}

    \begin{tabular}{lc}
        \toprule
        Componente & Versión \\
        \midrule
        Python & 3.12.10 \\
        NumPy & 2.4.6 \\
        pandas & 3.0.3 \\
        SciPy & 1.17.1 \\
        statsmodels & 0.14.6 \\
        scikit-learn & 1.9.0 \\
        seaborn & 0.13.2 \\
        \bottomrule
    \end{tabular}
\end{table}

Las dependencias del proyecto se registran adicionalmente en el archivo
\texttt{requirements.txt}. Esta separación entre código, datos,
resultados y dependencias reduce la dependencia del entorno local desde
el cual se ejecuta el análisis.

\subsection{Localización y carga del archivo}

El notebook no utiliza una ruta absoluta fija. En su lugar, busca la raíz
del repositorio recorriendo el directorio actual y sus directorios
superiores hasta encontrar la carpeta \texttt{.git} o el archivo
\texttt{README.md}. A partir de esa ubicación se define la ruta:

\begin{center}
    \texttt{data/german.data}
\end{center}

Si el archivo no se encuentra localmente, el procedimiento descarga el
paquete comprimido desde UCI, verifica que contenga
\texttt{german.data}, lo extrae y lo almacena en la carpeta
\texttt{data}. De esta manera, el notebook puede ejecutarse sin depender
de una copia manual ubicada fuera del proyecto.

El archivo original está separado por espacios y no contiene nombres de
columnas. Durante la lectura se asignan los veintiún nombres definidos
en el diccionario de variables: veinte predictores y la clase original.

Después de la carga se comprueba que la matriz tenga exactamente:

\[
1000\ \text{observaciones}
\quad \times \quad
21\ \text{columnas}.
\]

La ejecución se detiene mediante una excepción si las dimensiones no
coinciden con la documentación o si la variable
\variable{clase} contiene valores diferentes de \(1\) y \(2\).

\subsection{Conservación de los códigos originales}

Las variables categóricas se mantienen con sus códigos originales, por
ejemplo \texttt{A11}, \texttt{A34} o \texttt{A65}. Esta decisión conserva
la trazabilidad con la documentación de UCI y evita modificar la
información antes de definir el uso de cada predictor.

Para facilitar la lectura de tablas y visualizaciones se generan
etiquetas en español para cuatro variables especialmente relevantes:

\begin{itemize}
    \item propósito del crédito;
    \item estado de la cuenta corriente;
    \item nivel de ahorros;
    \item historial crediticio.
\end{itemize}

Las nuevas columnas utilizan el sufijo \texttt{\_lbl}. Por ejemplo,
\variable{proposito_lbl} contiene descripciones como ``Auto nuevo'',
``Educación'' o ``Negocio'', mientras que
\variable{proposito} conserva códigos como \texttt{A40}, \texttt{A46} o
\texttt{A49}.

Antes de construir cada etiqueta, el notebook compara los códigos
observados con las claves disponibles en el diccionario correspondiente.
La ejecución se detiene si aparece una categoría sin mapear. Esta
validación evita que un código inesperado se convierta silenciosamente
en un valor faltante.

\subsection{Verificaciones de integridad}

Las comprobaciones aplicadas a las variables originales se resumen en la
Tabla~\ref{tab:verificaciones-integridad}.

\begin{table}[H]
    \centering
    \caption{Verificaciones de integridad aplicadas al dataset.}
    \label{tab:verificaciones-integridad}

    \begin{tabularx}{0.94\textwidth}{
        @{}
        X
        >{\centering\arraybackslash}p{2.1cm}
        >{\centering\arraybackslash}p{1.8cm}
        @{}
    }
        \toprule
        Verificación & Criterio & Resultado \\
        \midrule

        Valores faltantes
        & Total igual a \(0\)
        & Cumple \\

        Filas completamente duplicadas
        & Total igual a \(0\)
        & Cumple \\

        Duración
        & Mayor que \(0\)
        & Cumple \\

        Monto
        & Mayor que \(0\)
        & Cumple \\

        Edad
        & Entre 18 y 100 años
        & Cumple \\

        Tasa de la cuota
        & Valores en \(\{1,2,3,4\}\)
        & Cumple \\

        Antigüedad en la residencia
        & Valores en \(\{1,2,3,4\}\)
        & Cumple \\

        Número de créditos en el banco
        & Valores en \(\{1,2,3,4\}\)
        & Cumple \\

        Número de dependientes
        & Valores en \(\{1,2\}\)
        & Cumple \\

        Clase original
        & Valores en \(\{1,2\}\)
        & Cumple \\

        Respuesta recodificada
        & Valores en \(\{0,1\}\)
        & Cumple \\

        \bottomrule
    \end{tabularx}
\end{table}

Las verificaciones no detectaron valores faltantes ni filas completamente
duplicadas. Asimismo, las variables numéricas y ordinales se encontraron
dentro de los rangos documentados. Por esta razón, no fue necesario
eliminar observaciones, imputar datos ni corregir categorías.

No se eliminaron valores extremos de duración o monto. Aunque algunas
observaciones se encuentran alejadas del centro de sus distribuciones,
corresponden a registros válidos del dataset y pueden representar
operaciones crediticias reales. Su influencia se examina mediante
estadística descriptiva, visualizaciones, transformaciones y
diagnósticos de los modelos, en lugar de descartarlas automáticamente.

\subsection{Recodificación de la variable respuesta}

El archivo original define:

\[
\variable{clase} =
\begin{cases}
1, & \text{buen riesgo},\\
2, & \text{mal riesgo}.
\end{cases}
\]

Para el análisis se construye la respuesta binaria:

\[
Y_i=\variable{mal_credito}_i =
\begin{cases}
0, & \text{si } \variable{clase}_i=1,\\[3pt]
1, & \text{si } \variable{clase}_i=2.
\end{cases}
\]

De esta forma, el mal crédito queda definido como el evento positivo. La
decisión es coherente con el objetivo del estudio y con la matriz de
costos: un falso negativo representa un crédito de mal riesgo clasificado
como bueno, es decir, el error al que la fuente asigna el mayor costo.

La recodificación también facilita la interpretación de las métricas.
El \textit{recall} de la clase positiva representa directamente la
proporción de malos créditos detectados correctamente por el modelo.

Además, se crea \variable{clase_lbl} como una variable categórica
ordenada con dos etiquetas:

\begin{center}
    \textit{Buen riesgo}
    \qquad y \qquad
    \textit{Mal riesgo}.
\end{center}

Esta variable se utiliza únicamente para mejorar la lectura de tablas y
figuras; los cálculos se realizan con \variable{mal_credito}.

La recodificación produce \(700\) observaciones de buen riesgo y \(300\)
de mal riesgo, equivalentes al \(70\,\%\) y \(30\,\%\) del conjunto,
respectivamente. Esta distribución será considerada posteriormente al
separar los datos y evaluar las métricas del clasificador.

\subsection{Tratamiento de categorías informativas}

Los códigos \texttt{A65} y \texttt{A124} requieren una consideración
especial. El primero representa ``nivel de ahorros desconocido o sin
cuenta'', mientras que el segundo representa ``propiedad desconocida o
inexistente''. Ambos son códigos definidos por la fuente y no valores
faltantes generados durante la carga.

Por tanto, se mantienen como categorías válidas. Reemplazarlos por
valores nulos eliminaría una señal potencialmente relevante del proceso
de evaluación crediticia. Sin embargo, su interpretación es limitada,
porque cada categoría combina dos situaciones diferentes. Por ejemplo,
no es posible distinguir entre la ausencia de una cuenta de ahorros y la
falta de información sobre ella.

De manera similar, \variable{tasa_cuota} y
\variable{anios_residencia} se conservan como escalas ordinales
codificadas entre 1 y 4. Estos números representan niveles ordenados, no
porcentajes exactos ni años observados directamente. Cuando se utilicen
como predictores lineales, esa simplificación deberá considerarse al
interpretar sus coeficientes.

\subsection{Transformación del monto}

El análisis exploratorio muestra que el monto presenta una distribución
marcadamente asimétrica hacia la derecha. Después de comprobar que todos
sus valores son positivos, se construye:

\[
\variable{log_monto}_i
=
\log\left(\variable{monto}_i\right).
\]

La transformación logarítmica reduce la influencia relativa de los
montos más altos, comprime la escala y permite representar relaciones
multiplicativas de una forma más cercana a la linealidad. No sustituye al
monto original: ambas variables se conservan y se utilizan según el
objetivo de cada análisis.

La pertinencia de esta transformación se evalúa posteriormente mediante
el histograma, el gráfico Q--Q y el ajuste por máxima verosimilitud de
distribuciones lognormal y gamma.

\subsection{Separación entre preprocesamiento y modelamiento}

Las operaciones descritas hasta este punto se realizan sobre el conjunto
completo porque corresponden a validaciones estructurales, recodificaciones
determinísticas y creación de etiquetas.

La separación entre entrenamiento y prueba no forma parte de esta etapa.
Se introduce posteriormente, antes del ajuste de la regresión logística,
mediante una partición estratificada de \(70\,\%\) para entrenamiento y
\(30\,\%\) para prueba. La estratificación conserva la proporción de mal
crédito en ambos subconjuntos.

Las variables categóricas tampoco se reemplazan globalmente por códigos
numéricos arbitrarios. En los modelos de \texttt{statsmodels} se
incorporan mediante términos categóricos explícitos, dejando una
categoría de referencia. De esta manera, la representación utilizada por
cada modelo queda vinculada a su especificación y puede interpretarse de
forma transparente.

\subsection{Resultado del preprocesamiento}

El resultado es una base con las \(1000\) observaciones originales, sin
imputaciones ni exclusiones, acompañada por etiquetas legibles y tres
variables derivadas principales:

\begin{itemize}
    \item \variable{mal_credito}, respuesta binaria;
    \item \variable{clase_lbl}, etiqueta de presentación;
    \item \variable{log_monto}, transformación logarítmica del monto.
\end{itemize}

Este conjunto constituye la entrada común para el análisis exploratorio,
la inferencia frecuentista, los modelos de regresión y el análisis
bayesiano desarrollados en las secciones siguientes.

% Contenido previsto:
% - Carga reproducible.
% - Validación de dimensiones, tipos, faltantes y duplicados.
% - Recodificación de variables.
% - Separación entrenamiento/prueba.
% - Pipeline de transformación y prevención de fuga de información.
% - Semilla y entorno computacional.
